\chapter[DONNÉES DE TÉLÉMÉTRIE]{Données de télémétrie}\label{donnuxe9es-de-tuxe9luxe9muxe9trie}

DTI \texttt{T}. Le champ Information contient la télémétrie.

\section{Format Telemetry Report}\label{format-telemetry-report}

\begin{itemize}
\tightlist
\item
  Sequence Number sur 3 caractères --- typiquement 3 chiffres, ou les 3
  lettres \texttt{MIC} (dans ce cas, virgule optionnelle avant la
  première valeur analogique).
\item
  5 valeurs analogiques 8 bits non signées (000-255, 3 chiffres
  décimaux).
\item
  1 valeur digitale 8 bits (8 caractères \texttt{0} ou \texttt{1}).
\end{itemize}

Format utilisé par le Kantronics KPC-3+ et l'APRS MIM (Micro Interface
Module).

\begin{center}
\begin{tabular}{|c|c|c|c|c|c|c|c|c|}
\hline
\textbf{Champ} & \lit{T} & \lit{\#}xxx, & aaa, & aaa, & aaa, & aaa, & aaa, & bbbbbbbb \\
 & DTI & Seq. No & Analog 1 & Analog 2 & Analog 3 & Analog 4 & Analog 5 & Digital \\
\hline
\textbf{Octets} & 1 & 5 & 4 & 4 & 4 & 4 & 4 & 8 \\
\hline
\end{tabular}
\end{center}
\noindent + Comment (n octets, optionnel)

Exemples : - \texttt{T\#005,199,000,255,073,123,01101001} -
\texttt{T\#MIC199,000,255,073,123,01101001}

\section{Définition on-air des
paramètres}\label{duxe9finition-on-air-des-paramuxe8tres}

Les données brutes sont interprétables librement. En pratique, un
utilisateur APRS peut définir les paramètres (coefficients quadratiques,
sens des bits binaires) à tout moment, en envoyant ces définitions comme
messages APRS. Les autres stations qui reçoivent ces messages savent
alors comment interpréter.

Quatre messages :

\begin{itemize}
\tightlist
\item
  \textbf{Parameter Name} (\texttt{PARM.})
\item
  \textbf{Unit/Label} (\texttt{UNIT.})
\item
  \textbf{Equation Coefficients} (\texttt{EQNS.})
\item
  \textbf{Bit Sense / Project Name} (\texttt{BITS.})
\end{itemize}

L'adressee de ces messages est l'indicatif de la station qui émet la
télémétrie. Ex : N0QBF lance un ballon \texttt{N0QBF-11} → les 4
messages sont adressés à \texttt{N0QBF-11}.

Voir chapitre 14 pour le format des messages.

\section{Parameter Name Message}\label{parameter-name-message}

Contient les noms (N) des 5 canaux analogiques et des 8 canaux digitaux.

\begin{center}
\small
\begin{tabular}{|c|c|c|c|c|c|c|c|c|c|c|c|c|c|}
\hline
\textbf{Champ} & PARM. & A1 & \lit{,}A2 & \lit{,}A3 & \lit{,}A4 & \lit{,}A5 & \lit{,}B1 & \lit{,}B2 & \lit{,}B3 & \lit{,}B4 & \lit{,}B5 & \lit{,}B6 & \lit{,}B7 \\
\hline
\textbf{Octets} & 5 & 1-7 & 1-7 & 1-6 & 1-6 & 1-5 & 1-6 & 1-5 & 1-4 & 1-4 & 1-4 & 1-3 & 1-3 \\
\hline
\end{tabular}
\end{center}
\noindent + \texttt{,B8} (1-3 octets). Décomptes incluant les virgules séparatrices.

\textbf{Note} : les largeurs de champs varient (héritage historique des
limitations d'écran). Les décomptes d'octets incluent les virgules
séparatrices. La liste peut s'arrêter à n'importe quel champ.

Ex :
\texttt{:N0QBF-11\ :PARM.Battery,Btemp,ATemp,Pres,Alt,Camra,Chut,Sun,10m,ATV}

\section{Unit/Label Message}\label{unitlabel-message}

Unités des analogiques (U) et labels des digitaux (L). Mêmes largeurs de
champs que \texttt{PARM.}.

Ex :
\texttt{:N0QBF-11\ :UNIT.v/100,deg.F,deg.F,Mbar,Kft,Click,OPEN,on,on,hi}

\section{Equation Coefficients
Message}\label{equation-coefficients-message}

Trois coefficients (a, b, c) pour chacun des 5 canaux analogiques.
Formule appliquée à la valeur brute \texttt{v} :

\begin{verbatim}
value = a × v² + b × v + c
\end{verbatim}

\begin{center}
\begin{tabular}{|c|c|c|c|c|c|}
\hline
\textbf{Champ} & \texttt{EQNS.} & A1 (\texttt{a,b,c}) & ,A2 & ,A3 & ,A4 \\
\hline
\end{tabular}
\end{center}
\noindent + \texttt{,A5}. Chaque canal analogique = 3 coefficients (a,b,c). La liste peut s'arrêter à n'importe quel champ.

La liste peut s'arrêter à n'importe quel champ.

Ex :
\texttt{:N0QBF-11\ :EQNS.0,5.2,0,0,.53,-32,3,4.39,49,-32,3,18,1,2,3}

Application : canal A1 = tension batterie en 1/100 volts, coefficients
\texttt{a=0,\ b=5.2,\ c=0}. Si valeur brute \texttt{v=199} :

\begin{verbatim}
voltage = 0 × 199² + 5.2 × 199 + 0
        = 1034.8 (1/100 volts)
        = 10.348 V
\end{verbatim}

\section{Bit Sense / Project Name
Message}\label{bit-sense-project-name-message}

\begin{itemize}
\tightlist
\item
  Motif 8 bits (0 ou 1) indiquant le sens de chaque canal digital
  correspondant au label.
\item
  Nom du projet associé à la station.
\end{itemize}

\begin{verbatim}
BITS. | b1b2b3b4b5b6b7b8 | Project Title (0-23)
\end{verbatim}

Ex :
\texttt{:N0QBF-11\ :BITS.10110000,N0QBF\textquotesingle{}s\ Big\ Balloon}

Interprétation : si B1=1, la caméra a cliqué. Si B2=0, le parachute
s'est ouvert. Etc.

